\section{Related Work}
\label{sec:related}

In this section, we review the existing literature relevant to our work, categorizing them into two streams: sequential recommendation and generative recommendation.

\vspace{+3pt}
\textbf{Sequential Recommendation.} Sequential recommendation focuses on predicting a user's next interaction based on their historical behavior sequences. Early methods relied on convolutional neural networks (CNNs)~\cite{Caser} or recurrent neural networks (RNNs)~\cite{GRU4Rec} to capture sequential dependencies. Motivated by advances in natural language processing, Transformer-based models such as SASRec~\cite{SASRec} and BERT4Rec~\cite{BERT4Rec} have become widely adopted, leveraging self-attention to model long-range dependencies efficiently. $\text{S}^3$-Rec~\cite{S3-Rec} further improves performance by maximizing mutual information across multiple types of contextual data. Despite their success, these methods generally follow a pure ID-based paradigm, representing items as randomly initialized atomic identifiers. This design inherently lacks semantic awareness, which restricts knowledge transfer across items, reduces generalization to cold-start scenarios, and limits scalability to large item collections.

\vspace{+3pt}
\textbf{Generative Recommendation.} Generative recommendation has been largely inspired by advances in LLMs~\cite{qwen3, GeneRec,hou2025generative,kuaishou_survey,li-etal-2024-large}. Early methods indexed items using purely text-based identifiers~\cite{P5,IDGenRec,BIGRec}. For example, BIGRec~\cite{BIGRec} represents items using their titles and fine-tunes LLMs for recommendation under a bi-step grounding paradigm. BinLLM~\cite{BinLLM} transforms collaborative embeddings obtained from external models into binary sequences, thereby aligning item representations with formats readily consumable by LLMs. While still compatible with LLM inputs, these approaches struggled to capture item-level collaborative signals~\cite{CoLLM,BinLLM,llara}. Subsequent approaches reformulate recommendation as a sequence-to-sequence task by tokenizing item sequences, providing a framework for jointly modeling semantic relevance and collaborative patterns~\cite{EAGER,SEATER,TIGER,LETTER,DiscRec,LC-Rec}. For example, TIGER~\cite{TIGER} represents each item with a Semantic ID—a tuple of codewords—and uses a Transformer to predict the next one. LETTER~\cite{LETTER} further introduces collaborative and diversity regularization to enhance the quality of the learned identifiers. OneRec~\cite{Onerec-v2} employs RQ-KMeans for tokenization, successfully replacing cascaded pipelines in industrial applications. 

However, most existing methods rely on a rigid two-stage pipeline in which the tokenizer is frozen during recommender training. This disjoint optimization renders the identifiers agnostic to downstream objectives, resulting in suboptimal representations~\cite{BLOGER}. More recent work, such as ETEG-Rec~\cite{ETEGRec}, explores alternating optimization to address this issue but still lacks a unified high-level objective. In contrast, our UniGRec addresses this critical gap in generative recommendation by leveraging soft identifiers, allowing the recommendation loss to jointly supervise both models. Combined with carefully designed training strategies, this enables fully end-to-end alignment between the tokenizer and recommender, bridging the disconnect present in prior two-stage pipelines.